\documentclass{article}
\usepackage{fullpage}
\usepackage{parskip}
\usepackage{standalone}
\usepackage{graphicx}
\usepackage{booktabs}
\usepackage{subcaption}
\usepackage{hyperref}
\usepackage{xurl}
\usepackage{amsmath}
\usepackage{amsfonts}
\usepackage{amsthm}
\usepackage{empheq}
\usepackage{xfrac}
\usepackage{enumitem}
\usepackage[table]{xcolor}
\usepackage{xfrac}
\usepackage{bbm}
\usepackage{siunitx}
\usepackage{array}
\usepackage{tikz}
\usetikzlibrary{arrows}
\usetikzlibrary{decorations.pathreplacing}
\usetikzlibrary{calc}
\usetikzlibrary{shapes.geometric}
\usetikzlibrary{decorations.markings}
\usetikzlibrary{shapes.geometric}
\usetikzlibrary{patterns,snakes}
\usetikzlibrary{decorations.text}
\usetikzlibrary{arrows.meta}
\usepackage{cases}
\usepackage{tcolorbox}
\usepackage{arydshln}

\newtheorem{definition}{Definition}
\newtheorem{remark}{Remark}

\title{\bf Modelling Deadlock Resolution in Open Restricted Queueing Networks}
\author{Geraint I. Palmer-Liyu}

\begin{document}
\maketitle

\section{Introduction}\label{sec:introduction}

\section{Preliminaries and System Model}\label{sec:recap}
This work builds directly on the work of \cite{palmer2018modelling}, which
considered deadlock detection in discrete event simulations (DES) of open
restricted queueing networks, and analysed Markovian models of deadlocking
queueing system.
The simulation software used is the Ciw library in Python \cite{palmer2019ciw},
although the methodology is appropriate for general DES.
We are concerned with open restricted queueing networks that experience Type I
blocking.
Queueing networks are described as open if customers can enter and leave the 
system from the exterior, while restricted networks contain at least one service
centre with limited queueing space or capacity before it, turning away
externally arriving customers when full, and inducing blocking for internally
routed customers when full.

This paper considers Type I blocking only \cite{onvural1986equivalencies}: once
a customer finishes their service and wishes to join another node's queue, they
are blocked from joining that queue if it is full; while blocked that customer
remains with its server until capacity becomes available at its destination;
during this time that server is unavailable to begin another customer's service
This gives rise to the idea that one customer can be blocked by another
customer: that is if customer A at node 1 is blocked from entering node 2
because it is full, then customer A is waiting on the customers who are with the
servers at node 2 to exit, hence customer A is blocked by the customers who are
with the servers of node 2.
\textit{Deadlock occurs when a subset of blocked customers are blocked by
customers in that subset only.}

In \cite{palmer2018modelling} a graph theoretic approach to detecting deadlock
during a simulation run is given:
Let $D(t) = (V, E(t))$ be the state digraph of a queueing system, with set of
vertices $V$ corresponding to the servers of all service centres, and a set of
directed edges $E(t)$ representing blockage relationships: if a customer at the
$k$th server of node $i$ is blocked from entering node $j$, then there are
directed edges from the vertex corresponding to node $i$'s $k$th server to every
vertex corresponding to the servers of node $j$.
Then it is shown that \textit{deadlock occurs at time $t$ if and only if $D(t)$
contains a knot}, that is a strongly connected component with no sink.
Figure~\ref{fig:exampledigraph} shows an example of a three node queueing
network in deadlock with the corresponding state digraph.

\begin{figure}
\begin{center}
\includestandalone[width=0.75\textwidth]{../images/example_digraph}
\end{center}
\caption{Example of a three node network in deadlock, and the state digraph
$D(t)$ that describes the blocking relationships. Customers with servers $A_1$,
$B_2$ and $C_1$ are blocked to node $B$, while the customer with server $B_1$ is
blocked to node $A$. The knot within $D(t)$ is highlighted - note that server
$C_1$ is not part of the knot, as no deadlocked customer is waiting for this
server to be free.}
\label{fig:exampledigraph}
\end{figure}

This paper builds on this graph theoretic foundation to consider deadlock
resolutions with no perceptible disruption to the customers' or the servers'
experiences.
First let's consider the unblocking behaviours, that is the cascading sequence
of servers freeing customers and capacity freeing up that occurs when a customer
blocking other customers is able to free a server.
By understanding this behaviour we then design an deadlock resolution method in
Section~\ref{sec:resolution}.


\section{Sequences of Unblockages}\label{sec:unblockages}
At any time $t$ during the simulation, some subset of the customers may be
blocked.
Define a blockage event as one specific instance of a specific customer being
blocked, and write $B_k$ as the $k$th blockage to occur during the simulation
run.
Now:

\begin{definition}
Let $(B_n)_{n \in \mathcal{N}(t)}$ be the sequence of active blockage events at
time $t$, where $\mathcal{N}(t) \subset \mathbb{N}$, and $|\mathcal{N}|$ is
bounded above by the number of servers in the queueing network.
\end{definition}

Blockage events are associated with both individual customers and the servers
who are occupied with the blocked customers.
Recall that $V$, the set of vertices of the state digraph $D(t)$, represents the
set of all servers, therefore there is a mapping between blockage events and
servers:

\begin{definition}
Let $v: (B_n)_{n \in \mathcal{N}(t)} \rightarrow V$ be an injective mapping
between blockage events and servers, such that $v(B_i) \in V$ represents the
server associated with blockage event $B_i$.
\end{definition}

Note that $v(\cdot)$ is injective as each blockage event must be associated with
a server, and each server can only be with one customer at a time, and so can
only be associated with one blockage event at a time.
$v(\cdot)$ is not surjective, as not all servers necessarily have to contain
blocked customers.

Let's consider how the system evolves over time $t$.
Both the set of active blockages $\mathcal{N}(t)$ and the sequence
$(B_n)_{n \in \mathcal{N}(t)}$ evolve with $t$.
Let's also introduce $M(t)$, the total number of blockages that have occurred so
far in the simulation run.
Then with each blockage and each unblockage the system state is updated
simultaneously by:

\begin{itemize}
  \item When a customer at server $v^{\star}$ is blocked to node $i$ with
  servers $V_i = \{v_j \in V \; | \; v_j \text{ is a server of node } i\}$:
    \begin{itemize}
      \item The total blockage counter $M(t)$ increments by $1$,
      \item The active blockage set $\mathcal{N}(t)$ expands to
      $\mathcal{N}(t) \cup \{M(t) + 1\}$,
      \item Create $B_{M(t) + 1}$ such that $v(B_{M(t) + 1}) = v^{\star}$,
      \item The set of edges of the state digraph $E(t)$ includes the new
      blockage relationships to become
      $E(t) \cup \{(v^{\star}, v_j) \; | \; v_j \in V_i\}$.
    \end{itemize}
  \item When the blockage event $B_k$ becomes unblocked:
    \begin{itemize}
      \item the index leaves the active blockage set $\mathcal{N}(t)$ to
      become $\mathcal{N}(t) \setminus \{k\}$,
      \item The associated edges are removed from the state digraph, so
      $E(t)$ becomes $E(t) \setminus \{(v(B_k), v_j) \; | \; v_j \in V\}$.
    \end{itemize}
\end{itemize}

Table~\ref{tbl:blockage_sequence} shows how a sequence of blockages and
unblockages can lead to the deadlock seen in Figure~\ref{fig:exampledigraph},
and how $(B_n)_{n\in\mathcal{N}(t)}$, $\mathcal{N}(t)$, $M(t)$, and $D(t)$
evolve at each step.
In this example there are three nodes: Node A with one server $A_1$, Node B with
two servers $B_1$ and $B_2$, and Node C with two servers $C_1$ and $C_2$.

\begin{table}
\begin{center}
\begin{tabular}{m{0.38\textwidth} m{0.08\textwidth} m{0.1\textwidth} m{0.14\textwidth} m{0.13\textwidth} }
\toprule
\vspace{0.5mm}\large{Event}\vspace{0.5mm} & \vspace{0.5mm}\large{$M(t)$}\vspace{0.5mm} & \vspace{0.5mm}\large{$\mathcal{N}(t)$}\vspace{0.5mm} & \vspace{0.5mm}\large{$(B_n)_{n\in\mathcal{N}(t)}$}\vspace{0.5mm} & \vspace{0.5mm}\large{$D(t)$}\vspace{0.5mm} \\
\midrule
Customer at $B_2$ is blocked to Node B & 1 & $\{1\}$          & $B_1$                & \vspace{1mm}\includestandalone[width=0.12\textwidth]{../images/digraph_step1} \\ \midrule
Customer at $A_1$ is blocked to Node B & 2 & $\{1, 2\}$       & $B_1, B_2$           & \vspace{1mm}\includestandalone[width=0.12\textwidth]{../images/digraph_step2} \\ \midrule
Customer at $C_2$ is blocked to Node C & 3 & $\{1, 2, 3\}$    & $B_1, B_2, B_3$      & \vspace{1mm}\includestandalone[width=0.12\textwidth]{../images/digraph_step3} \\ \midrule
Customer at $C_1$ is blocked to Node B & 4 & $\{1, 2, 3, 4\}$ & $B_1, B_2, B_3, B_4$ & \vspace{1mm}\includestandalone[width=0.12\textwidth]{../images/digraph_step4} \\ \midrule
Customer at $B_2$ is unblocked         & 4 & $\{2, 3, 4\}$    & $B_2, B_3, B_4$      & \vspace{1mm}\includestandalone[width=0.12\textwidth]{../images/digraph_step5} \\ \midrule
Customer at $B_1$ is blocked to Node A & 5 & $\{2, 3, 4, 5\}$ & $B_2, B_3, B_4, B_5$ & \vspace{1mm}\includestandalone[width=0.12\textwidth]{../images/digraph_step6} \\ \midrule
Customer at $C_2$ is unblocked         & 5 & $\{2, 4, 5\}$    & $B_2, B_4, B_5$      & \vspace{1mm}\includestandalone[width=0.12\textwidth]{../images/digraph_step7} \\ \midrule
Customer at $B_2$ is blocked to Node B & 6 & $\{2, 4, 5, 6\}$ & $B_2, B_4, B_5, B_6$ & \vspace{1mm}\includestandalone[width=0.12\textwidth]{../images/digraph_step8} \\
\bottomrule
\end{tabular}
\caption{An example sequence of blockages with unblockages.}
\label{tbl:blockage_sequence}
\end{center}
\end{table}

An unblockage can itself trigger further unblockages, for example if there
exists a set of blockages $\{B_i, B_j, ...\}$ which are all blocked by $v(B_k)$,
then unblocking $B_k$ will allow one of $\{B_i, B_j, ...\}$ to become unblocked.
This choice is governed by two conditions given in
Remark~\ref{rmk:unblocking_sequence}.

\begin{remark}\label{rmk:unblocking_sequence}
Unblocking $B_k$ will also trigger the unblocking of $B_h$ if:
\begin{enumerate}
  \item \label{pnt:necessity} $(v(B_h), v(B_k)) \in E(t)$, that is that the edge
  between the server associated with $B_h$ and the server associated with $B_k$
  in the state digraph; or in other words, that $B_h$ is blocked by $B_k$; and
  \item \label{pnt:choice}
  $h = \min(i \in \mathcal{N}(t) \; | \; (v(B_i), v(B_k)) \in E(t))$, that is
  out of all the active blockages caused by $B_k$, $B_h$ was the first to become
  blocked.
\end{enumerate}
\end{remark}

Note that condition~\ref{pnt:necessity} is a necessity, an unblockage can only
occur when there is spare capacity at the destination node, which is triggered
by the original unblockage.
Condition~\ref{pnt:choice} meanwhile is a choice: it would be natural in a FIFO
service discipline that those customers who have waited longest should be
unblocked first, retaining the eventual queueing order at the destination node;
though in other service disciplines other selection criteria may be more
appropriate.
As an example, Figure~\ref{fig:example_unblockages} shows a system with a
sequence of blockages, and Table~\ref{tbl:unblockage_sequence} shows the
sequence of unblockages that would occur \textit{instantaneously} when server
$B_1$ releases their customer and frees capacity at Node B.
This sequence of cascading unblockages eventually unblocks servers $C_1$, $B_2$
and $A_1$, while server $C_2$ remains blocked.

To illustrate the arbitrariness of condition~\ref{pnt:choice} above, consider
the case where server $A_1$ was the first to get blocked.
Once server $B_1$ releases their customer and creates capacity at Node B, $A_1$
is unblocked, but as there are no other blockages caused by $A_1$,
condition~\ref{pnt:necessity} is not satisfied, and so there are no more
unblockages.
Hence servers $B_2$, $C_1$ and $C_2$ all remain blocked.

\begin{figure}
\begin{center}
\includestandalone[width=0.675\textwidth]{../images/example_unblockages}
\end{center}
\caption{Example of a sequence of blockages and corresponding state digraph
$D(t)$, for which Table~\ref{tbl:unblockage_sequence} shows the sequence of
unblockages when the customer at server $B_2$ exits the system.}
\label{fig:example_unblockages}
\end{figure}

\begin{table}
\begin{center}
\begin{tabular}{m{0.38\textwidth} m{0.08\textwidth} m{0.1\textwidth} m{0.14\textwidth} m{0.13\textwidth} }
\toprule
\vspace{0.5mm}\large{Event}\vspace{0.5mm} & \vspace{0.5mm}\large{$M(t)$}\vspace{0.5mm} & \vspace{0.5mm}\large{$\mathcal{N}(t)$}\vspace{0.5mm} & \vspace{0.5mm}\large{$(B_n)_{n\in\mathcal{N}(t)}$}\vspace{0.5mm} & \vspace{0.5mm}\large{$D(t)$}\vspace{0.5mm} \\
\midrule
Customer at $B_1$ is released            & 4 & $\{1, 2, 3, 4\}$ & $B_1, B_2, B_3, B_4$ & \vspace{1mm}\includestandalone[width=0.12\textwidth]{../images/unblock_step1} \\ \midrule
Customer at $C_1$ is unblocked to Node B & 4 & $\{1, 2, 4\}$    & $B_1, B_2, B_4$      & \vspace{1mm}\includestandalone[width=0.12\textwidth]{../images/unblock_step2} \\ \midrule
Customer at $B_2$ is unblocked to Node C & 4 & $\{2, 4\}$       & $B_2, B_4$           & \vspace{1mm}\includestandalone[width=0.12\textwidth]{../images/unblock_step3} \\ \midrule
Customer at $A_1$ is unblocked to Node B & 4 & $\{2\}$          & $B_2$                & \vspace{1mm}\includestandalone[width=0.12\textwidth]{../images/unblock_step4} \\
\bottomrule
\end{tabular}
\caption{An example sequence of cascading unblockages.}
\label{tbl:unblockage_sequence}
\end{center}
\end{table}

\begin{tcolorbox}[colback=orange!5!white,colframe=orange!75!black,title=TODO]
Show that this is the true protocol for Ciw and SIMUL8. Do this (in an Appendix)
by designing a deterministic system where the sequence of unblockages does
indeed change macro behaviour: show we get the same results in Ciw as we do in
SIMUL8.
\end{tcolorbox}

\section{Deadlock Resolution}\label{sec:resolution}

\subsection{Induced Sequence of Unblockages}
For a cascading sequence of unblockages, which `untangles' the blockages built
up over time, to occur, it is necessary for the customer at the head of the
sequence to be blocked by servers who have the potential to create capacity at
their node.
That is, blocked by servers that are themselves are not blocked, or that the
system is not deadlocked.

We introduce the idea of artificially \textit{inducing} a sequence of
unblockages by creating dummy capacity, triggering a cascading sequence of
unblockages.
The aim is that this will clear the deadlocked state seamlessly, that is without
breaking the queueing capacity of Type 1 blocking rules for any meaningful
duration of simulation time.

\begin{definition}
At deadlock, a sequence of unblockages is \textsc{induced} by temporarily
creating extra capacity for one customer at one of the nodes.

The induced sequence of unblockages is \textsc{feasible} if by the end of the
sequence that extra capacity is not in use, and can therefore be removed.
\end{definition}

As sequences of unblockages happen instantaneously with respect to simulation
time, a feasible induced sequence of unblockages will not cause any customer to
experience any non-zero duration of time at the dummy capacity, and the nodes
themselves will not experience any non-zero duration of time with extra
capacity.
Therefore this feasible induced sequence of unblockages plays no role in the
modelled system, and is purely a simulation mechanic that stands as a proxy for
mutual communication between deadlocked service centres.

\begin{remark}\label{rmk:not_all_feasible}
Not all induced subsequences are feasible under the conditions set out in
Remark~\ref{rmk:unblocking_sequence}.
\end{remark}

The claim in Remark~\ref{rmk:not_all_feasible} can be shown with a simple
counter-example: consider the three node queueing network in deadlock given in
Figure~\ref{fig:counter_example}, and consider inducing a sequence of
unblockages by adding dummy capacity to Node A.
Then customer at $B_1$ will become unblocked, filling the dummy capacity and
leaving capacity at Node B; which in turn unblocks the customer at $C_1$,
filling the capacity at Node B and leaving capacity at Node C.
However no customer was originally blocked by $C_1$, and so no more unblockages
can occur, and the dummy capacity introduced at Node A is still used.
Hence this is not a feasible induced sequence of unblockages.

\begin{figure}
\begin{center}
\includestandalone[width=0.75\textwidth]{../images/counter_example}
\end{center}
\caption{Example of a deadlock where the induced sequence of unblockages by
adding dummy capacity at Node A would not be feasible under the conditions set
out in Remark~\ref{rmk:unblocking_sequence}.}
\label{fig:counter_example}
\end{figure}


\section{Markov Models of Deadlock Resolution}\label{sec:markov}

\bibliographystyle{plain}
\bibliography{refs}

\end{document}